\section{Lezione 9} \hfill 12/03/2026
\subsection{Castagno - Castanea}
È una delle principali specie forestali italiane (dopo l'abete rosso e faggio); si può trovare nella regione avanalpica e nel piano collinare.\\
È una specie che è stata per molto tempo considerata esotica (es. Grecia, Turchia), ma si è visto, mediante analisi polliniche, che è sempre stato presente in Italia.\\
Il suo uso è dovuto ai Greci, che hanno iniziato ad usarlo; successivamente i romani se ne sono appropriati, non tanto per l'alimentazione, ma come per materiale strategico bellico (da costruzione); per questo motivo lo hanno diffuso in tutta l'area dell'impero (es. Spagna, Portogallo, Romania), fino all'Inghilterra meridionale.\\
Nel medioevo si è iniziato a mangiare le castagne, che sono diventate un alimento fondamentale per le popolazioni collinari: quindi, c'è stata una ulteriore espansione del castagno per motivi alimentari.\\
Il castagno è una specie sporadica: in natura non forma popolamenti puri. Quando ci si trova davanti ad un castagneto, si sa che l'uomo ha influenzato quel territorio.\\
Un castagneto, se abbandonato, dopo una cinquantina d'anni tende a diventare quello che era prima (es. querceto misto).\\
Il castagno è una specie:
\begin{itemize}
    \item acidofila (vuole substrati acidi, non si sviluppa bene in substrati basici);
    \item fortemente eliofila;
    \item che vuole suoli fertili, ricchi e ben drenati (non sopporta ristagni idrici);
    \item molto longeva;
    \item ha una capacità pollonifera caulinare eterna (non ha limiti temporali al ricaccio);
    \item con una pasciona quasi annuale. La pasciona consiste nella produzione di molto seme, ciclica, che aumenta il successo della rinnovazione, perchè:
        \begin{enumerate}
            \item è uno sforzo di energia emettere ogni anno dei semi nuovi (es. una pianta che vive 200-300 anni non ha necessità di produrre seme ogni anno);
            \item limita il numero di predatori dei semi (es. cinghiali, scoiattoli), ma li emette solamente quando c'è possibilità di successo (se c'è molto seme, i predatori non fanno in tempo di moltiplicarsi a sufficienza e quindi alcuni semi riescono ad attecchire).
        \end{enumerate}
        La grande produzione di seme ha portato il castagno ad essere stato utilizzato molto.
    \item il suo areale di temperatura è tra gli 11 e 13 gradi e pluviometrico superiore a 800 mm (con almeno 100 mm nel trimestre estivo);
    \item la maturazione del seme è estiva.
\end{itemize}
Il castagno produce sia castagne che legno, ma (quasi) mai si ritraggono contemporaneamente dallo stesso popolamento.\\
\subsection{Castagneto da frutto}
Le castagne:
\begin{itemize}
    \item servono all'alimentazione come farina, caldarroste, castagne secche, alimentazione animale;
    \item hanno il difetto di non conservarsi o di conservarsi molto male;
    \item maturano tendenzialmente tutte assieme, tra inizio e fine ottobre;
    \item sono state selezionate delle varietà che si adattano meglio ai diversi usi (es. per farina, per caldarroste, per animali, da seccare, i marroni);
    \item alcune specie sono precoci, alcune sono più tardive.
\end{itemize}
Le castagne non sono molto digeribili.\\
Un riccio contiene tendenzialmente tre castagne, e quella centrale diventa poco piccola o veniva abortita.\\
Un marrone è una varietà più grande della castagna, dove vengono abortite le castagne laterali e quella centrale diventa molto grossa.\\
In passato si avevano molte varietà di castagne (più di 500): i popolamenti non erano monovarietali, ma composti da diverse specie per diversi usi.\\
La varietà più comune era la ``nostrana'' (specificatamente per ogni area).\\
I castagneti da frutto venivano tendenzialmente innestati. L'innesto veniva fatto su un selvaggione (popolamento multivarietale), più o meno ad altezza uomo (in modo che fosse comodo da lavorare), su fusti che avevano al massimo 10-12 cm DBH. L'innesto era tendenzialmente fatto a spacco o a corona, fatto con due-tre marze per portainnesto (per garantire almeno l'attecchimento di una).\\
In questo modo, i popolamenti erano nel piano inferiore dei selvaggioni, mentre nel piano superiore erano composti dalla varietà voluta\\
Nelle piante grosse è possibile notare se c'è stato l'innesto, se:
\begin{itemize}
    \item è presente una cicatrice lineare sul tronco;
    \item le cortecce e le dimensioni, di marza ed innesto, sono diverse (essendo geneticamente diverse).
\end{itemize}
Da giovane, la corteccia è liscia e grigiastra, dopo i 30-35 anni si screpola, diventando infine a spirale destrorsa.\\
In un castagneto da frutto si piantavano poche piante, in modo che le chiome fossero molto ampie (per produrre così molto frutto). Sotto al castagneto c'era solamente l'erba, per facilitare la raccolta delle castagne; solitamente alla base c'erano delle aiuole, per fermare il movimento delle castagne cadute. I ricci venivano bruciati per limitare la propagazione delle malattie del frutto.\\
Il castagno europeo (sativa) è una delle 7-8 specie di castagno al mondo, di cui 4 importanti per il frutti, le restanti sono arberelli-arbusti. \\
Il castagno, essendo una specie sporadica e non formando naturalmente popolamenti puri, è soggetto (quando è in purezza) ad essere attaccato da malattie. Tra la fine ed inizio del XIX-XX secolo, c'è stato il picco della coltivazione del castagno.\\
L'uomo ha provocato il passaggio di una delle malattie, che è il cancro corticale: è endemico sul castagno cinese, con il quale conviveva.\\
Esistono 7-8 specie di castagno, di cui quattro sono per i frutti:
\begin{multicols}{2}
\begin{itemize}
    \item castagno europeo (Castanea sativa);
    \item castagno americano (Castanea dentata);
    \item castagno cinese (Castanea mollissima);
    \item castagno giapponese (Castanea crenata)
\end{itemize}
\end{multicols}
Il castagno americano venne attaccato dal cancro corticale, importato e trasmesso dalla specie cinese.\\
Il cancro è arrivato in Europa nei prima anni '20, dall'America. Non è ancora arrivato nella bassa Inghilterra.\\
Il cancro ammazza facilmente le piante nate da seme e meno facilmente quelle nate da polloni. Mentre la specie americana aveva una mortalità elevata, per quello europeo la percentuale di morte era superiore per gli individui nati da seme.\\
I castagneti da frutto sono delle fustaie (e non dei cedui), ed il cancro ammazzava soprattutto la chioma, dove veniva prodotto il frutto.\\
In passato è stato ipotizzato di piantare il castagno cinese in Europa, che però non ha sapore; è perciò stato ibridato quello europeo con quello cinese, ed anche in questo caso le castagne non avevano un sapore appetibile.\\
I castagneti sono stati abbandonati, ed oggi avviene solamente la coltivazione per la produzione delle caldarroste.\\
Pochi anni fa, grazie alla moda del BIO, è tornata in auge la coltivazione del castagno, ed alcuni di questi popolamenti sono stati ripresi per mano.\\
Dove il suolo è pesante e umido, o dove è soggetto a forti variazioni del tasso di umidità, c'è un'altra malattia: il ``mal dell'inchiostro".\\
Mentre con il cancro corticale la pianta muore, ma ricaccia polloni, con il mal dell'inchiostro la pianta muore molto velocemente e non ricacciando polloni.\\
Il legno di castagno è:
\begin{itemize}
    \item usato per paleria, travi, mobili, pali telegrafici;
    \item molto bello (simile alla rovere, ma meno giallo);
    \item possiede moltissimi tannini (è un anti-crittogramo naturale):
        \begin{enumerate}
            \item ha una enorme durabilità, è adatto per gli usi esterni (es. pali da vigna);
            \item è molto contenuto sia nella corteccia che nel legno;
            \item veniva usato in passato per la concia delle pelli;
            \item essendo ignifugo, brucia male. Infatti il legno di castagno non si presta bene per la combustione (occorrerebbe lasciarlo all'acqua per 3-5 anni, affinchè si dilavino i tannini).
        \end{enumerate}
\end{itemize}
Il legno di castagno è prodotto tendenzialmente da un ceduo: in teoria potrebbe essere gestito come ceduo semplice, ovvero senza rilascio di matricine (data l'elevata capacità pollonifera e l'elevata produzione di seme, che parte da 12-15 anni).\\
Le matricine potrebbero servire a rimpiazzare le ceppaie che muoiono. In qualunque caso, i regolamenti forestali prescrivono la matricinatura, con minimo di 20-30 matricine (a seconda della Regione il numero può cambiare).\\
Anche se non è obbligatorio, è prescritta la matricinatura affinchè si mantenga la volontà politica del fatto che il bosco rimanga, anche dopo la ceduazione. Le matricine sono piante sfigate perchè: sono piante nate da seme (facilmente attaccabili dal cancro), vivono la loro vita dentro al bosco finchè non vengono ceduati i vicini (si scalda facilmente, si ...).\\
Un altro problema del castagno, se si vuole ottenere legname, è la cipollatura. Distacco degli anelli, di pochi cm o di tutta la circonferenza, che appare solitamente nel terzo mezzo del tronco. In molte specie la cipollatura appare dopo alcuni traumi; nel castagno e nell'abete bianco la cipollatura è genetica. È difficile da evitare, ma si possono ridurre: tenere le piante giovani (più invecchia e poù aumenta il rischio di..), mantenere gli accrescimenti costanti nel tempo, far crescere il castagno il più velocemente e tantissimo (in stazioni fertili). \\
Il viterbese è uno dei posti migliori per la coltivazione del castagno, grazie al terreno di origine vulcanica.\\
Iniziando da una fustaia di castagno, dopo il primo taglio si ottiene una fustaia.\\
Ad ogni ceduazione 8-10 perc delle ceppaie muore\\
Il turno ottimale del ceduo dipende dagli assortimenti che si vuole ottenere (in base al diametro). \\
5-7 \\
pali da vigna\\
25-30 anni per i travi...\\
La ceduazione oggi è semplice ed il turno normalmente è di 16 anni. I pali devono essere dritti e bisogna che non ci siano sciabolature. In passato veniva fatto uno sfollo sulla ceppaia nel terzo anno, e veniva rilasciato il 50 percento dei polloni affinchè crescessero dritti: la procedura non si fa più, causa il cancro. Ora la scelta si fa a 6 anni, quando i polloni non muoiono più.\\
Una ceppaia si evolve:
\begin{itemize}
    \item la prima ceduazione viene fatta con fusti intorno a 10-15 cm; si formano tanti polloni (anche 200-250 per ceppaia);
    \item i polloni mangiano tutti dallo stesso apparato radicale, e crescono tutti in maniera omogenea;
    \item la competizione tra i polloni, per lo spazio e luce, porta alla morte di alcuni polloni, facendone rimanere 4;
    \item questi polloni diventano grossi;
    \item dopo 16 anni, avviene il taglio dei polloni e la ceppaia è più grande di quella iniziale;
    \item i polloni ricacciano verso l'esterno, poichè c'è più luce;
    \item il problema è che ad un certo punto il castagno inizia a marcire e dentro si forma un buco. La pianta diventa un anello, ingrossandosi nuovamente, parte dell'anello potrebbe morire;
    \item può succedere che le ceppaie si dividano. le ceppaie hanno tutte lo stesso patrimonio genetico. 
\end{itemize} 
Al primo anno ci sono 300 polloni ad ettaro (?) ed intorno al sesto anno ci sono i polloni definitivi (che ci stanno). La competizione intra ceppaia si risolve in una decina d'anni.\\
Nel bosco ci sono diverse ceppaie, che creano competizione tra di loro, per tutta la vita del bosco. Il ceduo di castagno cresce tutto assieme, e per vincere la competizione le piante tendono a chiudere e coprire;\\
400-500 numero di ceppaie ad ettaro, anche fino a 800-1000. Molto spesso ci sono cedui con 100-200 ceppaie ad ettaro, che derivano da antichi cedui.\\
Se si volessero diametri grossi, occorre evitare le stazioni poco fertili, dobbiamo far crescere le piante il più libere possibili e tagliarle quando si raggiunge il diametro effettivo.\\
Un problema recente del castagno è il cinipide galigeno, che è arrivato una ventina di anni fa, da un'importazione abusiva di marze, ammazzando o riducendo fortemente gli incrementi. È stato importato il parassita